\documentclass[11pt,a4paper]{article}

\usepackage[utf8]{inputenc}
\usepackage[T1]{fontenc}
\usepackage[french]{babel}
\usepackage{lmodern}
\usepackage[a4paper,margin=2.0cm]{geometry}
\usepackage{microtype}
\usepackage{enumitem}
\usepackage{titlesec}
\usepackage{parskip}
\usepackage{graphicx}
\usepackage{xcolor}
\usepackage{listings}
\usepackage{booktabs}
\usepackage{tabularx}
\usepackage[hidelinks]{hyperref}

\titleformat{\section}
  {\normalfont\Large\bfseries}{\thesection.}{0.6em}{}
\titlespacing*{\section}{0pt}{0.9em}{0.25em}

\setlist[itemize]{leftmargin=*,topsep=2pt,itemsep=1pt,label=\textbullet}

\emergencystretch=3em
\tolerance=2000

\definecolor{usdabg}{RGB}{248,248,242}
\definecolor{usdakw}{RGB}{180,30,80}
\definecolor{usdacmt}{RGB}{120,120,120}
\definecolor{usdastr}{RGB}{40,120,40}

\lstdefinelanguage{usda}{
  morekeywords={def,Xform,Mesh,Material,Shader,Scope,over,prepend,references,uniform,token,float,int,double,double3,float3,color3f,point3f,rel,UsdPreviewSurface},
  morestring=[b]",
  morecomment=[l]{\#},
  sensitive=true,
}
\lstset{
  language=usda,
  basicstyle=\ttfamily\scriptsize,
  keywordstyle=\color{usdakw},
  stringstyle=\color{usdastr},
  commentstyle=\color{usdacmt},
  backgroundcolor=\color{usdabg},
  frame=single,
  framesep=4pt,
  framerule=0pt,
  rulecolor=\color{usdabg},
  numbers=left,
  numberstyle=\tiny\color{gray},
  numbersep=6pt,
  breaklines=true,
  showstringspaces=false,
  xleftmargin=12pt,
  belowskip=4pt,
  aboveskip=4pt,
}

\title{\vspace{-1.5em}\textbf{TP3 : Composition d'une forêt structurée avec OpenUSD}}
\author{ST2SNE, Efrei Paris 2025-2026 \quad Anil BRAUN}
\date{}

\begin{document}
\maketitle
\vspace{-2em}

\begin{center}
\small Dépôt GitHub : \href{https://github.com/Anilito1/forest-openusd-tp3}{\texttt{github.com/Anilito1/forest-openusd-tp3}}
\end{center}

\vspace{0.3em}
\noindent\includegraphics[width=\linewidth]{forest_game_view.png}
\vspace{-0.2em}
\begin{center}\small\textit{Forest.unity, Game View, après import du \texttt{Forest.usda} et matérialisation URP/Lit}\end{center}

\section{Présentation de la scène}

\textbf{Thème} : forêt tempérée stylisée, organisée autour d'un \emph{chemin sinusoïdal} qui relie trois écosystèmes distincts. L'intention de conception n'est pas une simple dispersion d'arbres, mais la composition d'un paysage \emph{lisible} où chaque zone porte une fonction visuelle et écologique propre.

\textbf{Zones} :

\begin{itemize}
  \item \textbf{Zone\_Dense} (forêt fermée, à gauche) -- canopée dominée par les pins et les chênes, sous-bois de buissons, fougères et mousses. Point focal : un chêne \emph{Patriarche} à l'échelle 1.85 placé au cœur de la zone.
  \item \textbf{Zone\_Clearing} (clairière, au centre) -- ouverture lumineuse bordée de bouleaux. Le sol porte 3 parterres de fleurs en clusters, un \emph{fairy ring} de 8 champignons, et 4 chênes en livrée d'automne dont la couleur de feuillage est modifiée par un \texttt{over} OpenUSD. Un nid d'oiseau est perché à 2,8 m sur un des arbres périphériques.
  \item \textbf{Zone\_Rocky} (zone rocailleuse, à droite) -- bouleaux dispersés, 4 grands rochers et 6 petits, un \emph{cairn} de 6 pierres empilées, un \emph{hollow log} (tronc tombé entouré de mousses et champignons), et un cluster de 3 arbres morts (\texttt{DeadTree}) pour la tension visuelle.
  \item \textbf{Path} -- chemin tracé par une courbe sinusoïdale qui traverse les trois zones ; 16 cailloux balisent le tracé. Un rayon d'exclusion de 1,8 m écarte les arbres du chemin.
\end{itemize}

La scène totalise \textbf{131 instances} composées par référence vers \textbf{14 assets} \texttt{.usda} distincts.

\section{Biodiversité}

\renewcommand{\arraystretch}{1.05}
\begin{tabularx}{\linewidth}{lXll}
\toprule
\textbf{Catégorie} & \textbf{Assets} & \textbf{Demandé} & \textbf{Livré} \\
\midrule
Arbres                 & \texttt{TreePine}, \texttt{TreeOak}, \texttt{TreeBirch}, \texttt{DeadTree} & 3 & \textbf{4} \\
Rochers                & \texttt{RockBig}, \texttt{RockSmall}                      & 2 & 2 \\
Végétation basse       & \texttt{Bush}, \texttt{Fern}, \texttt{Flower}             & 3 & 3 \\
Éléments du sol        & \texttt{Moss}, \texttt{FallenLog}, \texttt{Mushroom}      & 2 & \textbf{3} \\
Vie animale            & \texttt{BirdNest}                                          & 1 & 1 \\
Zones paysagères       & Dense, Clairière, Rocheuse \texttt{+} chemin               & 3 & \textbf{3+1} \\
\bottomrule
\end{tabularx}

\textbf{Logique de répartition} : chaque catégorie est plus probable dans une zone qu'une autre. Les pins et les chênes saturent la zone dense (densité 16+10 sur $\sim$300 m\textsuperscript{2}). La clairière concentre fleurs et fougères au sol et garde son centre dégagé. La zone rocheuse cumule rochers, mousses et bois mort. Un échantillonnage par \emph{rejection sampling} avec distance minimale par catégorie évite les chevauchements, et la distance au chemin agit comme exclusion supplémentaire pour les arbres.

\textbf{Strates d'âge} : 1 patriarche scale 1.85, 3 jeunes pousses scale 0.45, le reste centré sur 1.0 avec jitter 20--40\%. \textbf{Tilt organique} : \texttt{rotateZ} aléatoire $\pm$3$^{\circ}$ sur tous les arbres et buissons. \textbf{Overrides de matériau} : 4 chênes automne (palette orange/jaune), 2 bouleaux à écorce sombre, 2 buissons fleuris (rose et blanc crème).

\section{Structure OpenUSD}

\textbf{Organisation générale} : un fichier maître \texttt{Forest.usda} et un dossier \texttt{assets/} contenant les 14 \texttt{.usda} sources. Le maître ne contient \emph{aucun} \texttt{Mesh} inline : chaque objet de la scène est un \texttt{Xform} avec une référence vers son asset source.

\subsection*{Hiérarchie}
\begin{lstlisting}
def Xform "Forest"
{
    def Xform "Ground"        (prepend references = @./assets/Ground.usda@</Ground>)        { ... }
    def Xform "Path"          { def Xform "Stone_00" (prepend references = ...) { ... } ... }
    def Xform "Zone_Dense"    { def Xform "Patriarch_Oak" ... ; def Xform "Pine_07" ... ; ... }
    def Xform "Zone_Clearing" { def Xform "AutumnOak_00" ... ; def Xform "FairyRing_Mushroom_03" ... ; ... }
    def Xform "Zone_Rocky"    { def Xform "Cairn_Stone_04" ... ; def Xform "HollowLog" ... ; ... }
}
\end{lstlisting}

\subsection*{Instance par référence + variation locale}
\begin{lstlisting}
def Xform "Pine_07" (
    customData = { string species = "pine" ; string zone = "dense" }
    prepend references = @./assets/TreePine.usda@</TreePine>
)
{
    double3 xformOp:translate = (-6.95, 0, 3.53)
    float   xformOp:rotateY   = 321.18
    float   xformOp:rotateZ   = -1.4
    float3  xformOp:scale     = (1.08, 1.08, 1.08)
    uniform token[] xformOpOrder = ["xformOp:translate","xformOp:rotateY","xformOp:rotateZ","xformOp:scale"]
}
\end{lstlisting}

\subsection*{Override de matériau (chêne d'automne)}
\begin{lstlisting}
def Xform "AutumnOak_00" (
    prepend references = @./assets/TreeOak.usda@</TreeOak>
)
{
    over "Materials" {
        over "Mat_Foliage" {
            over "Surface" { color3f inputs:diffuseColor = (0.85, 0.45, 0.10) }
        }
    }
    double3 xformOp:translate = (3.0, 0, -3.5)
    ...
}
\end{lstlisting}

\textbf{Pourquoi des références plutôt que des copies} : un seul \texttt{TreePine.usda} est chargé en mémoire, ré-utilisé pour les 17 instances de pin. Modifier la géométrie ou la couleur de feuillage de tous les pins se fait en éditant un unique fichier source. Le fichier maître \texttt{Forest.usda} reste léger (68 ko, 1240 lignes) malgré 131 instances.

\section{Projet Unity}

\textbf{GitHub} : \href{https://github.com/Anilito1/forest-openusd-tp3}{\texttt{github.com/Anilito1/forest-openusd-tp3}}

\textbf{Scène à ouvrir} : \texttt{unity/ForestProject/Assets/Scenes/Forest.unity} (Unity 6.x, URP).

\textbf{Pipeline d'import} (deux scripts Editor fournis dans \texttt{Assets/Editor/}) :

\begin{enumerate}
  \item \texttt{ForestImporter.cs} -- menu \textbf{Tools $\rightarrow$ Forest $\rightarrow$ Import Forest.usda}. Appelle \texttt{Unity.Formats.USD.ImportHelpers.InitForOpen} puis \texttt{ImportSceneAsGameObject}, instancie la racine \texttt{Forest}, et déclenche \texttt{UsdAsset.Reload(true)} pour expanser la hiérarchie complète (132 GameObjects).
  \item \texttt{ForestMaterialize.cs} -- menu \textbf{Tools $\rightarrow$ Forest $\rightarrow$ Materialize}. Parcourt les 373 \texttt{MeshRenderer} de la scène et leur assigne un matériau \texttt{URP/Lit} parmi 18, choisi par mot-clé sur le chemin de hiérarchie (\texttt{trunk}, \texttt{foliage}, \texttt{autumnoak}, \texttt{cairn}, \texttt{mushroom/cap}, \ldots). Étape nécessaire car le SDK Unity USD ne convertit pas \texttt{UsdPreviewSurface} vers URP automatiquement.
\end{enumerate}

\textbf{Remarque sur les chemins Windows} : ouvrir le projet depuis un chemin \emph{sans accent ni espace}. Le SDK \texttt{com.unity.formats.usd} échoue sur des chemins contenant \texttt{é}, \texttt{à} ou des espaces ($\rightarrow$ \texttt{Cannot determine file format} ou \texttt{Failed to open layer}). Documenté dans le README.

\section{Difficultés rencontrées}

\textbf{1. Encodage des chemins Windows.} Le bug bloquant du TP. \texttt{InitForOpen} ne sait pas ouvrir un \texttt{.usda} dont le chemin absolu contient des caractères non-ASCII. Mon arborescence cours (\emph{Spécifications fonctionnelles\ldots}) déclenchait systématiquement l'erreur. Solution : projet Unity dans un chemin court et ASCII (\texttt{C:\textbackslash Users\textbackslash<user>\textbackslash Desktop\textbackslash ForestProject\textbackslash}).

\textbf{2. URP vs UsdPreviewSurface.} Le SDK USD Unity instancie correctement les meshes mais ne convertit pas les \texttt{UsdPreviewSurface} en matériaux URP. La scène apparaissait entièrement blanche avant matérialisation. Solution : script \texttt{ForestMaterialize.cs} qui crée 18 matériaux URP/Lit assignés par mot-clé sur le nom des GameObjects -- compromis acceptable, mais la couleur perd la précision millésimale du \texttt{diffuseColor} de chaque shader.

\textbf{3. Expansion de la hiérarchie.} \texttt{ImportSceneAsGameObject} renvoie un GameObject racine avec composant \texttt{UsdAsset} mais sans enfants visibles. L'appel explicite à \texttt{UsdAsset.Reload(true)} reconstruit la hiérarchie complète.

\textbf{Limites de la scène}. (i) Géométrie procédurale simple (cônes, cylindres, sphères) -- esthétique stylisée, pas photoréaliste. (ii) Le sol est un plan unique de 40\,m\,$\times$\,40\,m sans variations de relief. (iii) Pas d'animation : le format \texttt{.usda} le supporterait via \texttt{TimeSamples}, mais hors scope du TP.

\textbf{Améliorations possibles}. (a) Remplacer les meshes procéduraux par des assets Polyhaven importés en \texttt{.usdc}. (b) Préserver les \texttt{UsdPreviewSurface} via un \texttt{ScriptedImporter} URP custom. (c) Ajouter une variante \texttt{variantSet} \texttt{season} (\texttt{summer}/\texttt{autumn}/\texttt{winter}) sur chaque arbre au lieu de l'override ponctuel. (d) Couper le sol en plusieurs \texttt{.usda} référencés (\emph{payload}) pour permettre le chargement différé par zone.

\end{document}
